% Source: Weinberg GR, physical PDF pp. 84--87
%         (printed pp. 58--61).
% Coverage: Section 2.12, Representations of the Lorentz Group; equations
%           (2.12.1)--(2.12.19).
% Figures/tables/footnotes: optional-reading footnote; reference markers 11--12.
% Status: source-reviewed and compile-clean.
% Uncertainties: none.

\subsection{Representations of the Lorentz Group\optionalreading}
\label{sec:2.12}
\begingroup
\renewcommand\thefootnote{*}
\footnotetext{This section lies somewhat out of the book's main line of
development, and may be omitted in a first reading.}
\endgroup

The tensor formalism described in Section~\ref{sec:2.5} is perfectly adequate
for handling the problems of relativistic classical physics. However, there are
certain formal advantages in looking at these Lorentz transformation rules in a
more general way, from the perspective of the theory of the representations of
the homogeneous Lorentz group. We shall see in Section~12.5 that this approach
allows an elegant reformulation of the effects of gravitation on arbitrary
physical systems. Also, it is only in this way that we can deal with fields of
half-integer spin.

Under the general Lorentz transformation rule, a set of quantities \(\psi_n\)
transform under a Lorentz transformation \(\Lambda^\mu{}_\nu\) into the new
quantities
\begin{equation}
  \psi'_n=\sum_m[D(\Lambda)]_{nm}\psi_m .
  \tag{2.12.1}\label{eq:2.12.1}
\end{equation}
In order for a Lorentz transformation \(\Lambda_1\), followed by a Lorentz
transformation \(\Lambda_2\), to give the same result as the Lorentz
transformation \(\Lambda_1\Lambda_2\), it is necessary that the matrices
\(D(\Lambda)\) should furnish a \emph{representation} of the Lorentz group;
that is,
\begin{equation}
  D(\Lambda_1)D(\Lambda_2)=D(\Lambda_1\Lambda_2),
  \tag{2.12.2}\label{eq:2.12.2}
\end{equation}
with matrix multiplication now understood. For instance, if \(\psi\) is a
contravariant vector \(V^\mu\), then \(D(\Lambda)\) is simply
\begin{equation}
  [D(\Lambda)]^\mu{}_\nu=\Lambda^\mu{}_\nu ,
  \tag{2.12.3}\label{eq:2.12.3}
\end{equation}
whereas for a covariant tensor \(T_{\mu\nu}\), the corresponding \(D\) matrix
is
\begin{equation}
  [D(\Lambda)]_{\mu\nu}{}^{\rho\sigma}
  =\Lambda_\mu{}^\rho\Lambda_\nu{}^\sigma .
  \tag{2.12.4}\label{eq:2.12.4}
\end{equation}
It is easy to check that Eqs.~\eqref{eq:2.12.3} and \eqref{eq:2.12.4} do
satisfy the group multiplication rule \eqref{eq:2.12.2}. We can compile a
catalogue of all possible Lorentz transformation rules by constructing the most
general representation of the homogeneous Lorentz group.

In fact, the most general true representations of the homogeneous Lorentz
group are provided by the tensor representations, such as
Eqs.~\eqref{eq:2.12.3} and \eqref{eq:2.12.4}, so we might expect that all
quantities of physical interest should be tensors. However, there are additional
representations of the \emph{infinitesimal} Lorentz group, the \emph{spinor
representations}, that play an important role in relativistic quantum field
theory. The infinitesimal Lorentz group consists of Lorentz transformations
infinitesimally close to the identity; that is,
\begin{equation}
  \Lambda^\mu{}_\nu=\delta^\mu{}_\nu+\omega^\mu{}_\nu,
  \qquad \lvert\omega^\mu{}_\nu\rvert\ll1 .
  \tag{2.12.5}\label{eq:2.12.5}
\end{equation}
In order for this to satisfy the fundamental condition \eqref{eq:2.1.2} for a
Lorentz transformation, we must have
\[
  (\delta^\mu{}_\rho+\omega^\mu{}_\rho)
  (\delta^\nu{}_\sigma+\omega^\nu{}_\sigma)\eta_{\mu\nu}
  =\eta_{\rho\sigma},
\]
or, to first order in \(\omega\),
\begin{equation}
  \omega_{\rho\sigma}=-\omega_{\sigma\rho},
  \tag{2.12.6}\label{eq:2.12.6}
\end{equation}
with the indices on \(\omega\) of course lowered with \(\eta\):
\[
  \omega_{\rho\sigma}\equiv\eta_{\rho\mu}\omega^\mu{}_\sigma .
\]
For such a transformation, the matrix representation \(D(\Lambda)\) must be
infinitesimally close to the identity:
\begin{equation}
  D(1+\omega)=1+\frac12\omega^{\mu\nu}\sigma_{\mu\nu},
  \tag{2.12.7}\label{eq:2.12.7}
\end{equation}
where \(\sigma_{\mu\nu}\) are a fixed set of matrices, which by virtue of
Eq.~\eqref{eq:2.12.6} can always be chosen to be antisymmetric in \(\mu\)
and \(\nu\):
\begin{equation}
  \sigma_{\mu\nu}=-\sigma_{\nu\mu}.
  \tag{2.12.8}\label{eq:2.12.8}
\end{equation}
For instance, for the tensor representations \eqref{eq:2.12.3} and
\eqref{eq:2.12.4}, we have
\begin{equation}
  [\sigma_{\mu\nu}]^\rho{}_\sigma
  =\delta^\rho{}_\mu\eta_{\nu\sigma}
   -\delta^\rho{}_\nu\eta_{\mu\sigma},
  \tag{2.12.9}\label{eq:2.12.9}
\end{equation}
\begin{equation}
\begin{aligned}
  [\sigma_{\mu\nu}]_{\rho\sigma}{}^{\lambda\kappa}
  ={}&\eta_{\mu\rho}\delta^\lambda{}_\nu\delta^\kappa{}_\sigma
     -\eta_{\nu\rho}\delta^\lambda{}_\mu\delta^\kappa{}_\sigma\\
    &+\eta_{\mu\sigma}\delta^\kappa{}_\nu\delta^\lambda{}_\rho
     -\eta_{\nu\sigma}\delta^\kappa{}_\mu\delta^\lambda{}_\rho .
\end{aligned}
  \tag{2.12.10}\label{eq:2.12.10}
\end{equation}
The matrices \(\sigma_{\mu\nu}\) are not allowed to be just any set of
constant matrices, but must be constrained so that \(D(\Lambda)\) satisfies
the group multiplication rule \eqref{eq:2.12.2}. It is convenient first to
apply this rule to the product \(\Lambda(1+\omega)\Lambda^{-1}\):
\[
  D(\Lambda)D(1+\omega)D(\Lambda^{-1})
  =D(1+\Lambda\omega\Lambda^{-1}).
\]
To zero order in \(\omega\), this simply says that \(1=1\), whereas to first
order, we must equate the coefficients of \(\omega_{\mu\nu}\) on both sides:
\begin{equation}
  D(\Lambda)\sigma_{\mu\nu}D(\Lambda^{-1})
  =\sigma_{\rho\sigma}\Lambda^\rho{}_\mu\Lambda^\sigma{}_\nu .
  \tag{2.12.11}\label{eq:2.12.11}
\end{equation}
If we now set \(\Lambda=1+\omega\) (not necessarily with the same \(\omega\))
and \(\Lambda^{-1}=1-\omega\), then this will be satisfied to first order in
\(\omega\) provided that \(\sigma\) satisfies the commutation relations
\begin{equation}
  [\sigma_{\mu\nu},\sigma_{\rho\sigma}]
  =\eta_{\rho\nu}\sigma_{\mu\sigma}
   -\eta_{\rho\mu}\sigma_{\nu\sigma}
   +\eta_{\sigma\nu}\sigma_{\rho\mu}
   -\eta_{\sigma\mu}\sigma_{\rho\nu},
  \tag{2.12.12}\label{eq:2.12.12}
\end{equation}
with square brackets denoting the usual matrix commutator
\[
  [u,v]\equiv uv-vu .
\]

The reader can easily check that the matrices \eqref{eq:2.12.9} and
\eqref{eq:2.12.10} do satisfy\linebreak Eq.~\eqref{eq:2.12.12}. The problem of finding
the general representations of the infinitesimal homogeneous Lorentz group is
thus reduced to the problem of finding all matrices that satisfy the
commutation relations \eqref{eq:2.12.12}.

These commutation relations can be put in a somewhat more familiar form by
defining the matrices
\begin{equation}
\begin{array}{lll}
  a_1=\frac12(-\ii\sigma_{23}+\sigma_{10}),&
  a_2=\frac12(-\ii\sigma_{31}+\sigma_{20}),&
  a_3=\frac12(-\ii\sigma_{12}+\sigma_{30}),\\[2mm]
  b_1=\frac12(-\ii\sigma_{23}-\sigma_{10}),&
  b_2=\frac12(-\ii\sigma_{31}-\sigma_{20}),&
  b_3=\frac12(-\ii\sigma_{12}-\sigma_{30}).
\end{array}
  \tag{2.12.13}\label{eq:2.12.13}
\end{equation}
Equation \eqref{eq:2.12.12} then takes the form
\begin{equation}
  \mathbf a\times\mathbf a=\ii\mathbf a ,
  \tag{2.12.14}\label{eq:2.12.14}
\end{equation}
\begin{equation}
  \mathbf b\times\mathbf b=\ii\mathbf b ,
  \tag{2.12.15}\label{eq:2.12.15}
\end{equation}
\begin{equation}
  [a_i,b_j]=0 .
  \tag{2.12.16}\label{eq:2.12.16}
\end{equation}
Equations \eqref{eq:2.12.14}--\eqref{eq:2.12.16} are simply the commutation
relations for a pair of independent angular momentum matrices. The rules for
constructing such matrices can be found in any book on nonrelativistic quantum
mechanics.\textsuperscript{\hyperref[ch2-ref-11]{11}} In the most general case
\(\mathbf a\) and \(\mathbf b\) are a direct sum of ``irreducible''
components, each characterized by an integer or half-integer \(A\) or \(B\),
with
\begin{equation}
  \mathbf a^2=A(A+1),\qquad
  \mathbf b^2=B(B+1),
  \tag{2.12.17}\label{eq:2.12.17}
\end{equation}
and with dimensionality \(2A+1\) or \(2B+1\). Thus the most general objects
\(\psi_n\), which transform linearly under infinitesimal homogeneous Lorentz
transformations, can be decomposed into ``irreducible'' pieces, characterized
by a pair of integers and/or half-integers \((A,B)\), each piece having
\((2A+1)(2B+1)\) components.

\begingroup
\emergencystretch=1.5em
A straightforward calculation shows that the contravariant vector
representation \eqref{eq:2.12.9}, as well as its covariant counterpart, has
\(A=B=\tfrac12\). Any tensor representation, such as
\eqref{eq:2.12.10}, can be regarded as a direct product of vector
representations, so it consists only of irreducible components with \(A+B\)
an integer; for instance, the general second-rank tensor representation
\eqref{eq:2.12.10} consists of irreducible components with \((A,B)\) equal
to \((1,1),(1,0),(0,1)\), and \((0,0)\). The representations in which
\(A+B\) is a \emph{half-integer} are quite distinct from the tensors, and are
called \emph{spinor representations}. The most familiar example is the Dirac
electron field, which consists of components with \((A,B)\) equal to
\((\tfrac12,0)\) and \((0,\tfrac12)\).
\par
\endgroup

The transformation property of any object under ordinary spatial rotations is
determined by its behavior with respect to infinitesimal Lorentz transformations
\eqref{eq:2.12.5} for which \(\omega_{i0}=0\), and hence by the structure of
the purely spatial components \(\sigma_{12},\sigma_{23},\sigma_{31}\) of
\(\sigma_{\mu\nu}\). From these components, we can construct a matrix vector
\begin{equation}
  \mathbf s=\mathbf a+\mathbf b
  =-\ii\{\sigma_{23},\sigma_{31},\sigma_{12}\}
  \tag{2.12.18}\label{eq:2.12.18}
\end{equation}
that according to Eqs.~\eqref{eq:2.12.14}--\eqref{eq:2.12.16} has the
commutation relations of an angular momentum:
\begin{equation}
  \mathbf s\times\mathbf s=\ii\mathbf s .
  \tag{2.12.19}\label{eq:2.12.19}
\end{equation}
\begingroup
\emergencystretch=1.5em
Any irreducible representation \((A,B)\) of the homogeneous Lorentz group can
be decomposed into pieces\textsuperscript{\hyperref[ch2-ref-11]{11}} with
\(\mathbf s^2\) equal to \(s(s+1)\), where \(s\) is an integer or
half-integer between \(\lvert A-B\rvert\) and \(A+B\); each term describes
excitations (e.g., particles) of spin \(s\). It follows then from
Eq.~\eqref{eq:2.12.18} that the tensor representations can describe only
excitations with integer spin, whereas the spinor representations describe only
excitations with half-integer spin.
\par
\endgroup

Finite Lorentz transformations can be built up by multiplying together an
infinite number of infinitesimal Lorentz transformations. In the same way, the
tensor representations of the infinitesimal Lorentz group can be used to
construct the tensor representations, such as Eqs.~\eqref{eq:2.12.3} and
\eqref{eq:2.12.4}, of the group of finite Lorentz transformations. However,
if we try to construct spinor representations of the finite Lorentz
transformations, we find that we can only get ``representations up to a
sign'';\textsuperscript{\hyperref[ch2-ref-12]{12}} that is, the group
multiplication law \eqref{eq:2.12.2} will occasionally have a minus sign on
the right-hand side. For instance, the product of two successive
\(180^\circ\) rotations about a given axis does not give the unit matrix, but
minus the unit matrix. The appearance of these minus signs means that a spinor
field itself cannot be a physical observable, though even functions of spinor
fields can be observables.
